\section*{Overview}

Coordinate compression is not an optimization trick in isolation. It is the preprocessing step that makes an
array-based solution legal once the original values are too large, negative, or sparse.

\section*{Problem-Driven Motivation}

Suppose an array contains values up to \(10^9\), possibly negative, and the intended solution is:

\begin{itemize}
  \item Fenwick tree over value ranks,
  \item segment tree over event coordinates,
  \item frequency array over distinct values.
\end{itemize}

The algorithmic idea is fine, but allocating arrays indexed directly by raw coordinates is impossible. What really
matters is usually not the exact coordinate size. It is the relative order.

That is exactly the situation where compression turns a dead idea into a working one.

\section*{Recognition Pattern}

Compression is usually the right move when:

\begin{itemize}
  \item values are huge but only \(O(n)\) of them ever appear,
  \item comparisons such as \(<\), \(\le\), or equality are the only important operations,
  \item the next structure wants small integer indices,
  \item the full set of referenced values is known offline.
\end{itemize}

Typical contest phrases are: ``values up to \(10^9\)'', ``negative coordinates'', ``sparse points'', and ``after sorting,
the actual distances do not matter.''

\section*{Derivation}

If an algorithm only depends on order, then replacing every value by its rank does not change the answer.

Formally, after sorting the distinct values
\[
v_0 < v_1 < \cdots < v_{k-1},
\]
map every original value \(x\) to the unique index \(id(x)\) such that \(v_{id(x)} = x\).

Then:

\begin{itemize}
  \item \(x < y \iff id(x) < id(y)\),
  \item \(x = y \iff id(x) = id(y)\).
\end{itemize}

That is enough for rank-based data structures.

What compression \emph{does not} preserve is geometric length. If an interval length or coordinate gap matters directly,
you must either keep the original values too, or transform the problem more carefully.

\section*{Worked Problem}

\paragraph{Problem.}
Count inversions in an array of length \(n \le 2 \cdot 10^5\), where values may be negative and as large as \(10^9\).

\paragraph{Why the naive approach fails.}
The naive double loop is \(O(n^2)\). A Fenwick tree over raw values is impossible because the value range is far too
large.

\paragraph{How compression fixes it.}

\begin{itemize}
  \item Collect all array values.
  \item Sort and deduplicate them.
  \item Replace each value by its rank.
  \item Process the array from left to right with a Fenwick tree over compressed ranks.
\end{itemize}

If the current value has compressed rank \(r\), then the number of earlier elements greater than it is:
\[
\texttt{seen\_so\_far} - \texttt{fenwick.prefix\_sum}(r+1).
\]

\paragraph{Why this is correct.}
Compression preserves the order ``greater than'' exactly. So counting bigger ranks is the same as counting bigger
original values.

\section*{Implementation Reasoning}

The practical questions that matter are:

\begin{itemize}
  \item \textbf{When do I collect values?} Before building the mapping, gather every value that any later operation will
  reference. In offline problems this includes query values, endpoints, and sometimes transformed values like \(r+1\).
  \item \textbf{Do I need reverse mapping?} Keep the sorted unique array if later logic needs the original values back.
  \item \textbf{0-index or 1-index?} Compression itself is naturally 0-indexed, but Fenwick trees are often 1-indexed.
  Convert once and keep it consistent.
\end{itemize}

The code below provides a reusable compressor plus an inversion-counting application.

\section*{Correctness Intuition}

Compression only changes labels. It does not change the sorted order of equal or distinct values. Therefore any
algorithm whose logic depends solely on order or equality behaves identically after remapping.

\section*{Complexity Analysis}

For \(m\) collected values:

\begin{itemize}
  \item build compression: \(O(m \log m)\),
  \item one lookup with \texttt{lower\_bound}: \(O(\log m)\),
  \item memory: \(O(m)\).
\end{itemize}

\section*{Common Pitfalls}

\begin{itemize}
  \item Forgetting to include values that appear only in future queries.
  \item Compressing interval endpoints but forgetting that a half-open model may also need \(r+1\).
  \item Treating compressed indices as if their differences were real geometric distances.
  \item Mixing 0-indexed compressed values with a 1-indexed Fenwick tree.
\end{itemize}

\section*{Variants and Failure Modes}

\begin{itemize}
  \item Compress pairs or tuples lexicographically when one scalar is not enough.
  \item Compress both axes in offline geometry or rectangle-query problems.
  \item If updates introduce previously unseen coordinates online, static compression is no longer enough; use maps,
  balanced trees, or rebuild offline.
\end{itemize}

\section*{Practice Problems}

\begin{itemize}
  \item Inversion counting.
  \item Offline rectangle counting after compressing one or two axes.
  \item Dynamic frequency or order-statistics queries on large sparse values.
\end{itemize}

\section*{References}

\begin{itemize}
  \item \href{https://codeforces.com/catalog}{Codeforces Catalog}.
  \item \href{https://cp-algorithms.com/data_structures/segment_tree.html}{cp-algorithms: common segment-tree uses with compressed coordinates}.
  \item \href{https://usaco.guide/plat/sweep-line}{USACO Guide: Sweep-line and offline geometry modules}.
\end{itemize}
